% ---------------------------------------------------------------------------
% Numeric macros used by shared/body.tex. Every number below was re-verified
% this session (\AsOf{}) against a primary source -- committed code, a
% committed report file, a currently-passing test, a git commit, or (for one
% group) a live PyPI lookup -- not copied from research/research.md, which is
% where this paper's four topics originate (Tier 4, items 4.1-4.4) but which
% this project's own house rule treats as a pointer, not a source, after
% several prior errors found in it this session. Provenance is grouped by
% section, matching this paper's four numbered sections.
%
% Evidence grade is noted per group, using this project's established
% vocabulary (see research/papers/silent_failures and
% research/papers/hab_precision_baserate):
%   [ARTIFACT]       a committed report/data file or a currently-passing test
%                     reproduces the number on demand.
%   [ARTIFACT, code]  a constant read directly from committed source code
%                     (not a report), equally reproducible and equally fixed.
%   [DOCUMENTATION]   descriptive prose in a committed module docstring,
%                     README, or comment -- reproducible in principle by
%                     re-running what it describes, but not backed by a
%                     stored output file or a named timing script.
%   [DATED]           a live, single-run measurement recorded with its date
%                     in a module docstring, re-runnable but with no stored
%                     output committed. Cited as of that date.
%   [EXTERNAL]        a fact about a third-party package, checked live
%                     against that package's own published metadata rather
%                     than assumed from this repo's comments about it.
%
% A note at the end of this file documents every place a source read this
% session disagreed with research.md, with CLAUDE.md, with another file in
% this same repository, or with itself.
% ---------------------------------------------------------------------------

\newcommand{\AsOf}{2026-09-10}
\newcommand{\NSections}{4}

% =============================================================================
% Section 1 -- An MLflow dependency inverted, not added
% =============================================================================

% --- This repo's own pins [ARTIFACT] ---
% machine_learning/requirements.txt:22-25,58 (pandas, pyarrow, mlflow-skinny),
% read directly \AsOf{}.
\newcommand{\PandasPinned}{3.0.5}
\newcommand{\PyarrowPinned}{25.0.0}
\newcommand{\MlflowSkinnyPinned}{3.6.0}

% --- Never fails a training run [ARTIFACT, code] ---
% machine_learning/marine_ml/tracking.py's `_guard` context manager (every
% logging call site is wrapped in it) and `track()`'s own try/except around
% mlflow.start_run/end_run, both read directly \AsOf{}.
\newcommand{\HabRunMinutes}{40}

% --- The backend/ boundary [ARTIFACT] ---
% backend/tests/test_training_run_record.py::test_the_backend_never_imports_mlflow
% and ::test_mlflow_is_not_a_backend_dependency, re-run \AsOf{}: 10 passed
% (both files, test_training_run_record.py + test_download_gebco.py together).
\newcommand{\BackendMlflowOffenders}{0}

% --- What the full `mlflow` package pins, checked against PyPI's own
% release metadata rather than assumed from this repo's comment [EXTERNAL,
% verified \AsOf{}] ---
% https://pypi.org/pypi/mlflow/3.6.0/json (the release matching the
% mlflow-skinny version pinned here; upload_time 2025-11-07): requires_dist
% includes "pandas<3" and "pyarrow<23,>=4.0.0" verbatim -- confirming this
% repo's own comment (requirements.txt:49-51, tracking.py:17-19) is accurate
% for the specific release it was written against.
\newcommand{\MlflowMatchedVersion}{3.6.0}
\newcommand{\MlflowMatchedPandasCeiling}{3}
\newcommand{\MlflowMatchedPyarrowCeiling}{23}
% https://pypi.org/pypi/mlflow/json (latest release at time of writing,
% 3.16.0): requires_dist has already relaxed to "pandas<4" and
% "pyarrow<26,>=4.0.0" -- both satisfied by this repo's actual pins above
% with no downgrade required. See the note at the end of this file: the
% "would force a downgrade" justification is a property of mlflow==3.6.0
% specifically, not a stable fact about "mlflow", and needs re-checking at
% the next upgrade rather than being treated as settled.
\newcommand{\MlflowLatestVersion}{3.16.0}
\newcommand{\MlflowLatestPandasCeiling}{4}
\newcommand{\MlflowLatestPyarrowCeiling}{26}

% =============================================================================
% Section 2 -- Reproducing a 52-model audit in under a second per variable
% =============================================================================

% --- Audit scale [ARTIFACT] ---
% research/data/results/meta.json (variables: 13, horizons: 4,
% as_of: "2026-08-10", duration_seconds: 5229.8); research/data/sites.csv
% (25 lines = 1 header + 24 sites); research/data/results/loso.json
% (a JSON list, len() == 1248, each entry one (variable, horizon, site) fit);
% research/data/results/ablation_meta.json (duration_seconds: 2411.5).
% All four read directly \AsOf{}.
\newcommand{\NVariablesAudited}{13}
\newcommand{\NHorizonsAudited}{4}
\newcommand{\NSitesAudited}{24}
\newcommand{\NModelsAudited}{52}
\newcommand{\NHeldoutFits}{1{,}248}
\newcommand{\MainExperimentSeconds}{5229.8}
\newcommand{\AblationExperimentSeconds}{2411.5}

% --- The freezing mechanism itself [ARTIFACT, code] ---
% backend/forecasting/history.py:61 (_DISK_TTL = timedelta(hours=6)),
% :65 (_MEMORY_TTL = timedelta(minutes=15)), :236-250 (frozen_cache(), gated
% on MARISAI_FROZEN_HISTORY_CACHE=="1"); backend/scripts/run_paper_experiments.py:79
% (AS_OF = date(2026, 8, 10)); grep confirms the env var is set only by
% run_paper_experiments.py, run_seasonality_ablation.py and
% export_research_dataset.py, never by a serving path. All read directly \AsOf{}.
\newcommand{\DiskCacheTtlHours}{6}
\newcommand{\MemoryCacheTtlMinutes}{15}
\newcommand{\PinnedAsOfDate}{2026-08-10}

% --- The per-variable cache speedup [DOCUMENTATION -- prose only, no stored
% timing artifact or named benchmark script; see the note at the end of this
% file] ---
% Stated identically in three places: research/README.md's "Two things that
% make the experiments reproducible" section; backend/forecasting/history.py:243's
% own frozen_cache() docstring ("~700s per variable"); and
% backend/scripts/run_paper_experiments.py:25's module docstring ("~700s").
% No committed benchmark log gives "0.8s" a reproducible source beyond this
% repeated prose. Graded separately from the ARTIFACT-backed counts above.
\newcommand{\AuditSecondsUncached}{700}
\newcommand{\AuditSecondsCached}{0.8}
% Derived from the two macros above (700 / 0.8 = 875), not a new measurement.
\newcommand{\AuditSpeedupFactor}{875}
% backend/forecasting/history.py's module docstring: an uncached
% Copernicus point-series fetch costs 8-13s [DOCUMENTATION].
\newcommand{\PointFetchSecondsLow}{8}
\newcommand{\PointFetchSecondsHigh}{13}

% =============================================================================
% Section 3 -- Access-pattern-dependent performance in ARCO zarr stores
% =============================================================================

% --- The one measured direction [DATED, 2026-08-10; no stored artifact] ---
% docs/ml-notes.md lines 32-40 (Machine learning conventions section),
% "Measured 2026-08-10 on one year of global monthly physics coarsened to
% 0.25deg: geo-series 7.5s, time-series 89.0s." Read directly \AsOf{}.
% This is a single named-date measurement in prose, not a stored benchmark
% file; the reversal for the opposite access pattern (bounded area, many
% timesteps) is stated in the same source as a consequence of the two
% services' opposite chunking layout, not as a second timed measurement --
% see the note at the end of this file.
\newcommand{\ArcoMeasuredDate}{2026-08-10}
\newcommand{\ArcoGeoSeriesSeconds}{7.5}
\newcommand{\ArcoTimeSeriesSeconds}{89.0}
% Derived (89.0 / 7.5 = 11.867, rounded), not a new measurement.
\newcommand{\ArcoSlowdownFactor}{11.9}

% --- A second, separate hazard: the missing depth bound [DOCUMENTATION] ---
% docs/ml-notes.md lines 27-31: the `_my_` reanalysis products carry 50 depth
% levels; passing minimum_depth/maximum_depth to open_dataset turns a request
% that "did not finish in 15 minutes" into "~60s". Distinct mechanism from
% the service-selection finding above (a missing server-side filter, not a
% choice between two services) -- kept as a separate macro group
% deliberately, not folded into the ARCO pair.
\newcommand{\ReanalysisDepthLevels}{50}
\newcommand{\ReanalysisUnboundedMinutes}{15}
\newcommand{\ReanalysisBoundedSeconds}{60}

% =============================================================================
% Section 4 -- ERDDAP answers 404 while reloading a dataset
% =============================================================================

% --- Retry policy constants and classification [ARTIFACT, code +
% currently-passing tests] ---
% backend/forecasting/history.py:319-320 (_FETCH_ATTEMPTS = 3,
% _RETRY_BACKOFF_SECONDS = (2.0, 6.0)), :335-376 (is_retryable);
% backend/tests/test_forecasting.py::test_permanent_client_errors_are_not_retried,
% ::test_a_404_is_retried_exactly_once,
% ::test_the_404_budget_does_not_leak_into_other_client_errors -- re-run
% \AsOf{}: 8 passed.
\newcommand{\FetchAttemptsMax}{3}
\newcommand{\RetryBackoffFirstSeconds}{2.0}
\newcommand{\RetryBackoffSecondSeconds}{6.0}
\newcommand{\RetryAttemptsFourZeroFour}{2}
\newcommand{\RetrySecondsFourZeroFour}{2.0}
\newcommand{\RetryAttemptsFourHundred}{1}
\newcommand{\RetryAttemptsFiveZeroThree}{3}
\newcommand{\RetrySecondsFiveZeroThree}{8.0}
% Derived: 2.0 + 6.0 = 8.0, the total backoff a fully-retried (5xx) or
% wrongly-retried (permanent 4xx, pre-fix) request pays.
\newcommand{\RetryPenaltySeconds}{8}

% --- The corrected belief: GEBCO_2020 was never retired [ARTIFACT, git
% history + code comment] ---
% Commit 50438f01b30b717cb02137b57c61c86bac655b46, "fix(forecasting): retry
% an ERDDAP 404 exactly once", 2026-08-05 15:50 IST: commit message states
% the re-probe directly. Corroborated by
% backend/services/download/providers/gebco.py:27-35's own comment and
% docs/features/forecasting.md lines 72-80. All read/confirmed via `git log`
% and direct file reads \AsOf{}.
\newcommand{\GebcoReprobeDate}{2026-08-05}
\newcommand{\ErddapFlapSeconds}{100}

% --- The training-run cost of one transient failure [DOCUMENTATION, in-code
% comment, dated to commit 39df31409ccc5f8ea96e9325d67e84a0046a2e94
% (2026-08-05), no separate stored report] ---
% backend/forecasting/history.py:314-318's own comment: "A single ERDDAP 503
% during one training run cost 6 of 24 training points -- a quarter of the
% model's geography, including the entire Gulf Stream and Coral Sea".
\newcommand{\ErddapCostAffectedPoints}{6}
\newcommand{\ErddapCostTotalPoints}{24}

% --- The cost of retrying a permanent failure: two disagreeing sources
% [an internal inconsistency, both halves committed -- see the note at the
% end of this file] ---
% backend/forecasting/history.py:340-341 (is_retryable's own docstring) and
% backend/tests/test_forecasting.py:877-879
% (test_permanent_client_errors_are_not_retried's docstring) both say:
% "a warm 6s forecast/response" becomes 14.5s. docs/features/forecasting.md
% line 71 (and research.md's Tier-4 4.4 entry, which appears to be copied
% from it) instead say "a warm 2.2s response" becomes 14.5s. Both figures
% pair with the same +8s penalty (\RetryPenaltySeconds{}); only 6+8=14 is
% arithmetically close to 14.5, while 2.2+8=10.2 is not. \WarmForecastCodeSeconds{}
% is used as the primary figure in the prose on that basis; both are reported.
\newcommand{\WarmForecastCodeSeconds}{6}
\newcommand{\WarmForecastDocsSeconds}{2.2}
\newcommand{\WarmForecastAfterBadRetrySeconds}{14.5}

% ---------------------------------------------------------------------------
% Notes on discrepancies found while verifying this paper (\AsOf{})
% ---------------------------------------------------------------------------
%
% 1. [Section 1] research.md and this repo's own comments state, without a
%    version attached, that "the full package pins pandas<3 / pyarrow<23".
%    Checked live against PyPI: this is exactly correct for mlflow==3.6.0,
%    the release matching the mlflow-skinny version actually pinned here
%    (\MlflowMatchedPandasCeiling{}, \MlflowMatchedPyarrowCeiling{}) -- but
%    mlflow's latest release, 3.16.0, has already relaxed this to
%    pandas<\MlflowLatestPandasCeiling{} / pyarrow<\MlflowLatestPyarrowCeiling{},
%    both of which \PandasPinned{} and \PyarrowPinned{} already satisfy with
%    no downgrade required. The dependency-inversion rationale is accurate
%    as documented, but it is a fact about one pinned release, not a stable
%    fact about "mlflow" -- an upgrade of mlflow-skinny past 3.16.0 could
%    remove the need for the inversion, or a future mlflow release could
%    re-tighten it. Neither research.md nor this repo's own comments carry a
%    version number, which is what let this drift unnoticed.
%
% 2. [Section 1] marine_ml/tracking.py's own module docstring claims silent
%    non-tracking "is guarded by `active_run_id`, which the pipelines
%    print." `grep -rn "active_run_id" machine_learning/` (\AsOf{}) finds
%    exactly one occurrence: that docstring line itself. No pipeline prints
%    anything under that name; `RunHandle.run_id` exists but is never
%    surfaced to stdout under this or another label in the code read this
%    session. This paper does not rely on the claim.
%
% 3. [Section 2] research/README.md's own phrasing, "13 variables x 4
%    horizons x 24 sites = 52 models", reads as one product but is not one:
%    13 x 4 = 52 (\NModelsAudited{}) is the model count; the further x 24
%    (\NSitesAudited{}) gives \NHeldoutFits{} held-out fits (52 x 24 = 1248),
%    not a third factor of 52. This paper reports the two multiplications
%    separately rather than reproducing the compressed line.
%
% 4. [Section 2] The audit's *scale* (\NModelsAudited{} models,
%    \NHeldoutFits{} fits, \MainExperimentSeconds{}s / \AblationExperimentSeconds{}s
%    total run time) is ARTIFACT-graded, read directly from committed report
%    files. The *speedup* research.md and this paper both quote
%    (\AuditSecondsUncached{}s to \AuditSecondsCached{}s per variable) is
%    DOCUMENTATION-graded: it is stated consistently in three places but
%    traces to no committed benchmark log or named timing script, unlike the
%    counts beside it. The two claims in this section carry different
%    evidence weight and are graded separately rather than averaged together.
%
% 5. [Section 3] CLAUDE.md states that picking the wrong ARCO service "turns
%    a ~10s fetch into minutes/hours"; research.md's Tier-4 4.3 entry
%    tightens this to "hours" outright. Neither figure is supported by any
%    measurement found this session: the only timed comparison
%    (docs/ml-notes.md, \ArcoMeasuredDate{}) is \ArcoGeoSeriesSeconds{}s
%    against \ArcoTimeSeriesSeconds{}s -- an \ArcoSlowdownFactor{}x factor,
%    not hours -- and the separate depth-bound hazard's worst case is stated
%    as "did not finish in \ReanalysisUnboundedMinutes{} minutes", which is
%    a lower bound on how long it might have taken, not a measurement of
%    "hours". This paper reports only the two measured/stated pairs above
%    and does not repeat the "hours" framing. Separately, the task brief for
%    this paper states these numbers are "documented in the companion paper
%    research/papers/global_vs_regional/"; grepping that paper's shared/
%    directory for "arco-geo-series", "arco-time-series", "7.5s" and "89.0s"
%    found no matches. The only source located for these specific numbers is
%    docs/ml-notes.md; CLAUDE.md documents the same convention but with no
%    numbers attached.
%
% 6. [Section 4] See the macro group above ("two disagreeing sources"):
%    docs/features/forecasting.md and research.md both state the retrying-a-
%    permanent-failure cost example as a warm \WarmForecastDocsSeconds{}s
%    response becoming \WarmForecastAfterBadRetrySeconds{}s; the code comment
%    and test docstring that this exact scenario is regression-tested against
%    (backend/forecasting/history.py, backend/tests/test_forecasting.py) both
%    instead say \WarmForecastCodeSeconds{}s becoming
%    \WarmForecastAfterBadRetrySeconds{}s. Both pair with the same
%    \RetryPenaltySeconds{}s penalty; \WarmForecastCodeSeconds{} +
%    \RetryPenaltySeconds{} = 14, matching \WarmForecastAfterBadRetrySeconds{}
%    far more closely than \WarmForecastDocsSeconds{} +
%    \RetryPenaltySeconds{} = 10.2 does. This paper uses
%    \WarmForecastCodeSeconds{} as the reported figure on that arithmetic
%    ground, states both, and treats the mismatch as a second, independent
%    instance of documentation drift found while writing this paper -- the
%    same shape as note 1 above, in the same paper, about a different
%    section.
%
% 7. [Section 4] backend/tests/test_download_gebco.py's own docstring and
%    its test `test_the_retired_dataset_id_is_not_requested` still describe
%    NOAA CoastWatch's GEBCO_2020 as a dataset "the server retired" that "is
%    gone" -- the pre-\GebcoReprobeDate{} belief, not the corrected one
%    recorded in gebco.py's own comment and docs/features/forecasting.md
%    (the dataset was never gone; the server was flapping). The test itself
%    is not wrong -- it still correctly asserts the provider requests
%    Ifremer's `gebco2021` rather than `GEBCO_2020`, which remains true for
%    an unrelated reason (Ifremer stayed up throughout the flap) -- but its
%    prose was never updated to match the correction shipped in the same
%    area of the codebase. Confirmed still present and the full file still
%    passing (10/10, run together with test_training_run_record.py) \AsOf{}.
% ---------------------------------------------------------------------------
